\documentclass[conference]{IEEEtran}

\usepackage{amsmath}
\usepackage{booktabs}
\usepackage{graphicx}
\usepackage{url}

\title{Hardware-Assisted Online Softmax Merge for Attention State Reuse on Spatz}

\author{
  \IEEEauthorblockN{Wenxuan Tang}
  \IEEEauthorblockA{Local Spatz Prototype Study}
}

\begin{document}

\maketitle

\begin{abstract}
This draft reports the current Paper B prototype state: a cluster-local
online-softmax merge unit with a complete mixed-scalar datapath, fixed
256-segment exponential and reciprocal lookup approximations, and RTL
microbenchmark evidence. The current evaluation includes an attention-like
merge-chain fallback, but does not yet claim end-to-end attention speedup.
\end{abstract}

\section{Introduction}

Paper B extends the cluster-local merge-update platform from paper A toward a
complete online softmax merge engine for attention state reuse. The central
claim should be end-to-end or attention-like value, not only isolated
restricted-semantics microbenchmark speedup.

\section{Background}

This section should cover online softmax, attention tiling, state reuse, and
the Spatz/Snitch execution substrate~\cite{online-softmax2018,flashattention2022,spatz2023,snitch2020}.

\section{Architecture}

Describe the complete datapath:

\begin{align}
m_{new} &= \max(m_{old}, m_{tile}), \\
l_{new} &= l_{old}e^{m_{old}-m_{new}} +
           l_{tile}e^{m_{tile}-m_{new}}, \\
O_{new} &= O_{old}
          \frac{l_{old}e^{m_{old}-m_{new}}}{l_{new}} +
          O_{tile}
          \frac{l_{tile}e^{m_{tile}-m_{new}}}{l_{new}}.
\end{align}

This section should explicitly discuss `exp`, scaling, reciprocal/division,
pipeline latency, and TCDM scheduling.

\section{Numeric Design}

The current prototype uses a fixed approximation configuration. The
exponential path uses a Q1.23 lookup table with 256 linear segments over
$[-8,0]$, saturating below the range to zero and above the range to one. The
reciprocal path normalizes the input mantissa to $[1,2]$ and uses a 256-segment
Q1.23 lookup table with linear interpolation. No Newton refinement is used.

Against a full FP32 software reference on the synthetic full-merge cases, the
worst recorded maximum absolute error is $9.78\times10^{-6}$, the worst guarded
maximum relative error is $3.33\times10^{-6}$, and the worst mean absolute
error is $2.17\times10^{-6}$. This meets both the initial $10^{-3}$ target and
the exploratory $10^{-4}$ target for the current input set.

\section{Evaluation}

The current RTL benchmark validates five full mixed-scalar sweep points:

\begin{table}[h]
\centering
\caption{Full mixed-scalar SMU microbenchmark.}
\begin{tabular}{rrrrr}
\toprule
$N$ & $D$ & CPU cycles & SMU cycles & Speedup \\
\midrule
1 & 1 & 2327 & 1207 & 1.93$\times$ \\
4 & 8 & 17490 & 1478 & 11.83$\times$ \\
8 & 16 & 56103 & 2343 & 23.94$\times$ \\
8 & 32 & 97554 & 3371 & 28.94$\times$ \\
16 & 64 & 365620 & 9650 & 37.89$\times$ \\
\bottomrule
\end{tabular}
\end{table}

The CPU baseline in this table is an RTL-aligned fixed-point software reference
used for correctness gating and same-semantics comparison; it is not an
optimized vectorized software baseline. A synthetic attention-like fallback
with 8 rows, 4 merge blocks, and $D=32$ invokes the SMU path for each merge
step and records 388826 software-reference cycles versus 13432 SMU cycles
($28.95\times$), with 6188 TCDM accessed and 32 TCDM congested counter events.

\section{Limitations}

The current prototype only claims finite normal FP32 input coverage plus zero
for the online-softmax length fields where supported by the error policy. NaN,
Inf, subnormal, and both-zero length inputs are not part of the supported
numeric contract. The attention-like result is a merge-chain fallback rather
than a complete attention kernel; it excludes QK/PV computation, DRAM behavior,
runtime scheduling, and model-level accuracy. Resource, area, power, and
post-synthesis timing data are still missing.

\section{Conclusion}

The current data closes the minimal complete-datapath loop for the SMU
microbenchmark and provides attention-like merge-chain evidence. The next step
is parameter sweep and resource/timing evaluation, followed by integration into
a fuller attention runtime proxy.

\bibliographystyle{IEEEtran}
\bibliography{common/bib/refs}

\end{document}
